\chapter{Motiva\c tie}



\section{Contextul si problema}

Organizarea \textbf{conferințelor, meeting-urilor, concertelor} și a altor evenimente cu public și locuri limitate implică, în mod tradițional, procese fragmentate: promovarea prin canale disparate (rețele sociale, email, site-uri proprii), vânzarea biletelor sau înscrierilor în format fizic sau prin transfer bancar manual, și validarea participanților la intrare prin liste pe hârtie. Această abordare generează erori, pierderi de timp și dificultăți în reconcilierea plăților cu locurile rezervate.

Digitalizarea procesului — centralizarea informațiilor despre evenimente, automatizarea plăților și asocierea fiecărui bilet cu identitatea participantului — reprezintă un pas natural în contextul extinderii serviciilor online și al așteptărilor utilizatorilor privind comoditatea și securitatea tranzacțiilor \cite{fielding2000, stripe2024}.

Organizatorii (universități, firme, asociații culturale, săli de concert) nu dispun adesea de resurse pentru soluții enterprise costisitoare. În același timp, platformele internaționale generice nu sunt întotdeauna adaptate pieței locale: moneda RON, structura administrativă pe județe, sau necesitatea unui flux simplu de check-in administrat de organizator. Pentru \textbf{evenimente sportive} amatoriale, aceleași mecanisme (înscriere, plată, listă nominală) sunt utile, dar reprezintă un caz secundar față de conferințelor și concertelor.

Participanții au nevoie de o interfață clară: descoperire evenimente, filtrare, achiziție rapidă unuia sau mai multor bilete nominale, vizualizarea istoricului și gestionarea contului. Organizatorii au nevoie de un panou unic pentru publicare, monitorizare locuri disponibile și confirmare prezență la fața locului.




\section{Obiectivele lucrării}

Obiectivul principal este proiectarea și implementarea \textbf{aplicației web pentru organizarea și gestionarea evenimentelor}, care să îndeplinească următoarele obiective specifice:

\begin{enumerate}
	\item Implementarea unui modul de autentificare securizat: înregistrare cu verificare email, login JWT, resetare parolă, editare profil și ștergere cont.
	\item Oferirea unei interfețe publice pentru listarea evenimentelor viitoare, cu filtre după nume, județ și dată.
	\item Permiterea achiziționării biletelor nominale (unul sau mai multe nume per tranzacție) cu plată online prin Stripe Checkout, în moneda RON.
	\item Dezvoltarea unui panou de administrare pentru CRUD evenimente (inclusiv imagini), listă participanți și check-in.
	\item Gestionarea rolurilor utilizator/admin și promovarea/revocarea administratorilor cu parolă de autorizare.
	\item Documentarea arhitecturii, a modelului de date și a fluxurilor principale (UML, ERD).
\end{enumerate}

Lucrarea demonstrează integrarea tehnologiilor web moderne într-un produs coerent, testabil și prezentabil, cu respectarea unor practici de securitate recunoscute \cite{rfc7519, fowler2002}.


\section{Structura documentului}

Prezenta lucrare este organizată structural pe capitole, după cum urmează:

\vspace{0.2cm}
\begin{description}
	\setlength{\itemsep}{0.3cm}
	\item[Capitolul 2] Analizează în detaliu starea actuală a domeniului. Cuprinde studii relevante din literatură, o analiză a platformelor existente pe piață și o comparație directă cu soluția propusă.
	\item[Capitolul 3] Descrie în profunzime arhitectura sistemului și detaliile de implementare ale aplicației. Sunt incluse cerințele funcționale și non-funcționale, tehnologiile alese, diagramele UML, structura bazei de date și fragmente de cod sursă relevante.
	\item[Capitolul 4] Prezintă interfețele grafice ale aplicației prin intermediul unor capturi de ecran detaliate, mapate pe fluxurile principale de utilizare.
	\item[Capitolul 5] Sintetizează concluziile generale ale proiectului, evidențiază limitările identificate pe parcursul dezvoltării și propune viitoarele direcții de dezvoltare.
\end{description}


\section{Tipuri de evenimente vizate}

Aplicația web pentru organizarea și gestionarea evenimentelor este concepută în jurul unor formate frecvente în viața academică, culturală și profesională. Tabelul 1.1 sintetizează categoriile principale, actorii implicați și modul în care platforma le susține. 

\begin{table}[htbp]
	\centering
	\begin{tabular}{|l|l|p{7.5cm}|}
		\hline
		\textbf{Tip Eveniment} & \textbf{Actori Principali} & \textbf{Scenariu în Aplicație} \\ \hline
		Conferință / Meeting & Organizator, Participant & Publicare cu detalii complete (dată, locuri limitate, descriere, preț în RON), înregistrare online și check-in la intrare. \\ \hline
		Concert / Festival & Organizator, Public & Accent pe componenta vizuală (galerie de imagini), redirecționare rapidă către Stripe Checkout pentru achiziție bilet nominal. \\ \hline
		Competiție Sportivă & Admin, Sportiv & Înscriere pe bază de bilet nominal, check-in realizat de organizator la înregistrarea în concurs (caz secundar). \\ \hline
	\end{tabular}
	\caption{Tipuri de evenimente reflectate în aplicație}
	\label{tab:evenimente}
\end{table}

\subsection{Conferinte, meeting-uri si evenimente profesionale}

Universități, camere de comerț și firme de consultanță organizează periodic conferințe și meeting-uri cu număr limitat de locuri. Participanții trebuie să se înregistreze, să-și confirme identitatea (email verificat) și, uneori, să achite o taxă de participare. Platforma acoperă acest flux end-to-end: publicarea evenimentului cu dată, locație (județ, localitate) și capacitate, urmată de înregistrare online și emitere bilet nominal.

Pentru meeting-uri interne sau seminarii plătite, administratorul poate crea evenimente cu preț fix și număr redus de locuri; sistemul decrementează automat locurile disponibile la fiecare tranzacție confirmată. Check-in-ul din panoul admin permite validarea prezenței la recepție sau la intrarea în sală — util atât pentru conferințe, cât și pentru training-uri corporate.

\subsection{Concerte si evenimente culturale}

Concertele și festivalurile locale necesită o prezentare atractivă (titlu, descriere, mai multe imagini) și vânzare rapidă a biletelor. Aplicația permite încărcarea a până la opt imagini per eveniment, afișarea galeriei pe pagina de detalii și redirect către Stripe Checkout pentru plată securizată.

Filtrarea după județ și dată ajută publicul să descopere evenimente din apropiere — relevant pentru turnee naționale de concerte sau festivaluri itinerante. Biletele nominale asigură că fiecare loc este asociat unui participant, reducând revânzarea neautorizată și simplificând controlul la intrare.

\subsection{Evenimente sportive si utilizare complementara}

Modelul actual al aplicației (nume pe bilet, check-in administrat manual) se pretează și competițiilor sportive amatoriale sau turneelor locale, unde organizatorul trebuie să verifice identitatea participantului. Această nișă este secundară față de conferințe, concerte și meeting-uri: nu sunt implementate funcționalități specifice sportului (categorii de vârstă, echipe, program meciuri), dar fluxul general de înscriere și plată rămâne valabil.

Extinderi viitoare (categorii de bilete, abonamente de sezon) ar putea întări suportul sportiv fără a modifica arhitectura de bază.

\subsection{ Beneficii transversale}



Indiferent de tipul evenimentului, organizatorul beneficiaza de: 

	\begin{itemize}
		\item \textbf{Un singur panou} pentru toate formatele — nu este nevoie de soluții separate pentru conferințe vs. concerte.
		\item \textbf{Trasabilitate plăți} — fiecare bilet este legat de o sesiune Stripe și de un utilizator autentificat.
		\item \textbf{Adaptare locală} — județe românești, monedă RON, emailuri prin SMTP Brevo.
		\item \textbf{Control asupra datelor} — baza PostgreSQL self-hosted, fără dependență de un mar\-ketplace global.
	\end{itemize}
	
	\vspace{1em} % Spațiu între listă și următorul paragraf
	
	Această diversitate de cazuri de utilizare justifică alegerea unui model de date simplu (eveniment unic, preț unic, bilete nominale) care rămâne suficient de expresiv pentru majoritatea scenariilor întâlnite.
	
